\section{Experiments}
\begin{table}[t]
\centering
\small
\setlength{\tabcolsep}{3.0pt}
\caption{\small 
We report the small-world index (SWI $\uparrow$) of head-head functional graphs after applying SparseGPT and its variants at sparsity 0.5. 
 SWA (Head) and SWA (Layer) denote head-only and layer-only allocation strategies, respectively. 
\textbf{Higher SWI indicates better preservation} of small-world functional organization under pruning. Results are averaged over three seeds and reported as mean$_{\pm \text{std}}$.
The best result for each model is \colorbox{gray!15}{\textbf{bolded}}.
}
\resizebox{0.8\linewidth}{!}{

\begin{tabular}{lcccccc}
\toprule
\textbf{Model} 
& \textbf{SparseGPT} 
& \textbf{+ FARMS} 
& \textbf{+ ATP} 
& \textbf{+ SWA (Head)} 
& \textbf{+ SWA (Layer)} 
& \textbf{+ SWA} \\
\midrule
Qwen3 14B     & 0.1748$_{\scriptstyle \pm 0.06}$ & 0.1732$_{\scriptstyle \pm 0.06}$ & 0.1631$_{\scriptstyle \pm 0.09}$ & 0.1710$_{\scriptstyle \pm 0.09}$ & 0.1832$_{\scriptstyle \pm 0.08}$ & \colorbox{gray!15}{\textbf{0.1982}$_{\scriptstyle \pm 0.07}$} \\
Qwen3 8B      & 0.1470$_{\scriptstyle \pm 0.05}$ & 0.1560$_{\scriptstyle \pm 0.06}$ & 0.1418$_{\scriptstyle \pm 0.06}$ & 0.1545$_{\scriptstyle \pm 0.07}$ & 0.1606$_{\scriptstyle \pm 0.04}$ & \colorbox{gray!15}{\textbf{0.1775}$_{\scriptstyle \pm 0.05}$} \\
Qwen3 4B      & 0.1148$_{\scriptstyle \pm 0.10}$ & 0.1200$_{\scriptstyle \pm 0.09}$ & 0.1198$_{\scriptstyle \pm 0.05}$ & 0.1330$_{\scriptstyle \pm 0.08}$ & 0.1327$_{\scriptstyle \pm 0.05}$ & \colorbox{gray!15}{\textbf{0.1431}$_{\scriptstyle \pm 0.06}$} \\
\midrule
Llama3 8B    & 0.1285$_{\scriptstyle \pm 0.05}$ & 0.1420$_{\scriptstyle \pm 0.07}$ & 0.1227$_{\scriptstyle \pm 0.08}$ & 0.1341$_{\scriptstyle \pm 0.07}$ & 0.1420$_{\scriptstyle \pm 0.08}$ & \colorbox{gray!15}{\textbf{0.1599}$_{\scriptstyle \pm 0.06}$} \\
Llama3.2 3B  & 0.0772$_{\scriptstyle \pm 0.05}$ & 0.0880$_{\scriptstyle \pm 0.02}$ & 0.0898$_{\scriptstyle \pm 0.04}$ & 0.0877$_{\scriptstyle \pm 0.06}$ & 0.1005$_{\scriptstyle \pm 0.04}$ & \colorbox{gray!15}{\textbf{0.1126}$_{\scriptstyle \pm 0.04}$} \\
Llama3.2 1B  & 0.0558$_{\scriptstyle \pm 0.04}$ & 0.0520$_{\scriptstyle \pm 0.05}$ & 0.0524$_{\scriptstyle \pm 0.05}$ & 0.0559$_{\scriptstyle \pm 0.03}$ & 0.0661$_{\scriptstyle \pm 0.03}$ & \colorbox{gray!15}{\textbf{0.0700}$_{\scriptstyle \pm 0.04}$} \\
\bottomrule
\end{tabular}
}
\label{tab:swi_pruning_preservation}

\end{table}
\begin{table*}[t]
\centering
\scriptsize
\setlength{\tabcolsep}{2.0pt}
\caption{\small  Comparison of \textbf{WikiText Perplexity ($\downarrow$)}. We evaluate Base pruning methods (Wanda and SparseGPT) and their variants augmented with FARMS, ATP, and SWA at sparsity 0.5/0.7. Results are averaged over three seeds and reported as mean$_{\pm \text{std}}$. The best result in each setting is \colorbox{gray!15}{\textbf{bolded}}. SWA reduces perplexity and achieves the best values across all the models and sparsity levels.}
\begin{tabular}{ccc|llll|llll}
\toprule
\multirow{2}{*}{\textbf{Model}} & \multirow{2}{*}{\textbf{Size}} & \multirow{2}{*}{\textbf{Sparsity}} & \multicolumn{4}{c|}{\textbf{Wanda}} & \multicolumn{4}{c}{\textbf{SparseGPT}} \\
 &  &  & Base & + FARMS & + ATP & \textbf{+ SWA} & Base & + FARMS & + ATP & \textbf{+ SWA} \\
\midrule
\multirow{4}{*}{\shortstack{Llama 3.2}} & \multirow{2}{*}{1B} & 0.5 & 36.25$_{\scriptstyle \pm 0.18}$ & 33.90$_{\scriptstyle \pm 0.02}$ & 24.95$_{\scriptstyle \pm 0.04}$ & \cellcolor{gray!15}\textbf{22.42$_{\scriptstyle \pm 0.05}$} & 21.54$_{\scriptstyle \pm 0.29}$ & 21.14$_{\scriptstyle \pm 0.56}$ & 20.56$_{\scriptstyle \pm 0.34}$ & \cellcolor{gray!15}\textbf{20.47$_{\scriptstyle \pm 0.15}$} \\
 &  & 0.7 & 415.60$_{\scriptstyle \pm 32.1}$& 380.20$_{\scriptstyle \pm 20.1}$& 264.10$_{\scriptstyle \pm 21.7}$& \cellcolor{gray!15}\textbf{256.67$_{\scriptstyle \pm 18.00}$} & 157.20$_{\scriptstyle \pm 3.2}$& 136.80$_{\scriptstyle \pm 5.3}$& 106.60$_{\scriptstyle \pm 3.2}$& \cellcolor{gray!15}\textbf{105.00$_{\scriptstyle \pm 3.6}$} \\
 & \multirow{2}{*}{3B} & 0.5 & 19.71$_{\scriptstyle \pm 0.02}$ & 16.45$_{\scriptstyle \pm 0.09}$ & 14.42$_{\scriptstyle \pm 0.02}$ & \cellcolor{gray!15}\textbf{13.32$_{\scriptstyle \pm 0.12}$} & 23.98$_{\scriptstyle \pm 2.15}$ & 15.92$_{\scriptstyle \pm 1.85}$ & 13.53$_{\scriptstyle \pm 0.08}$ & \cellcolor{gray!15}\textbf{11.27$_{\scriptstyle \pm 0.98}$} \\
 &  & 0.7 & 177.80$_{\scriptstyle \pm 5.6}$ & 132.6$_{\scriptstyle \pm 5.3}$ & 125.1$_{\scriptstyle \pm 0.8}$ & \cellcolor{gray!15}\textbf{120.65$_{\scriptstyle \pm 2.31}$} & 70.70$_{\scriptstyle \pm 5.07}$ & 67.67$_{\scriptstyle \pm 2.93}$ & 51.51$_{\scriptstyle \pm 3.83}$ & \cellcolor{gray!15}\textbf{49.56$_{\scriptstyle \pm 0.92}$} \\
\multirow{2}{*}{\shortstack{Llama 3}} & \multirow{2}{*}{8B} & 0.5 & 11.86$_{\scriptstyle \pm 0.02}$ & 11.12$_{\scriptstyle \pm 0.03}$ & 11.50$_{\scriptstyle \pm 0.05}$ & \cellcolor{gray!15}\textbf{10.14$_{\scriptstyle \pm 0.03}$} & 11.69$_{\scriptstyle \pm 0.13}$ & 11.02$_{\scriptstyle \pm 0.87}$ & 10.72$_{\scriptstyle \pm 0.15}$ & \cellcolor{gray!15}\textbf{10.15$_{\scriptstyle \pm 0.08}$} \\
 &  & 0.7 & 140.50$_{\scriptstyle \pm 2.6}$ & 90.02$_{\scriptstyle \pm 3.00}$ & 85.93$_{\scriptstyle \pm 0.82}$ & \cellcolor{gray!15}\textbf{84.98$_{\scriptstyle \pm 0.20}$} & 55.94$_{\scriptstyle \pm 3.23}$ & 52.55$_{\scriptstyle \pm 1.98}$ & 40.40$_{\scriptstyle \pm 1.94}$ & \cellcolor{gray!15}\textbf{38.75$_{\scriptstyle \pm 1.29}$} \\
\midrule
\multirow{6}{*}{\shortstack{Qwen 3}} & \multirow{2}{*}{4B} & 0.5 & 21.91$_{\scriptstyle \pm 2.08}$ & 19.69$_{\scriptstyle \pm 0.16}$ & 18.88$_{\scriptstyle \pm 0.03}$ & \cellcolor{gray!15}\textbf{17.86$_{\scriptstyle \pm 0.49}$} & 20.52$_{\scriptstyle \pm 0.43}$ & 19.58$_{\scriptstyle \pm 0.41}$ & 20.44$_{\scriptstyle \pm 0.45}$ & \cellcolor{gray!15}\textbf{19.30$_{\scriptstyle \pm 0.77}$} \\
 &  & 0.7 & 101.80$_{\scriptstyle \pm 3.9}$ & 90.49$_{\scriptstyle \pm 4.84}$ & 64.78$_{\scriptstyle \pm 1.21}$ & \cellcolor{gray!15}\textbf{63.65$_{\scriptstyle \pm 0.27}$} & 50.93$_{\scriptstyle \pm 1.62}$ & 44.30$_{\scriptstyle \pm 2.26}$ & 46.45$_{\scriptstyle \pm 0.91}$ & \cellcolor{gray!15}\textbf{40.24$_{\scriptstyle \pm 0.38}$} \\
 & \multirow{2}{*}{8B} & 0.5 & 14.90$_{\scriptstyle \pm 0.48}$ & 13.20$_{\scriptstyle \pm 0.03}$ & 12.91$_{\scriptstyle \pm 0.06}$ & \cellcolor{gray!15}\textbf{11.90$_{\scriptstyle \pm 0.09}$} & 15.09$_{\scriptstyle \pm 0.07}$ & 12.61$_{\scriptstyle \pm 0.12}$ & 13.08$_{\scriptstyle \pm 0.09}$ & \cellcolor{gray!15}\textbf{12.10$_{\scriptstyle \pm 0.14}$} \\
 &  & 0.7 & 71.10$_{\scriptstyle \pm 1.86}$ & 65.46$_{\scriptstyle \pm 1.51}$ & 57.31$_{\scriptstyle \pm 1.28}$ & \cellcolor{gray!15}\textbf{54.05$_{\scriptstyle \pm 1.62}$} & 30.05$_{\scriptstyle \pm 0.66}$ & 25.13$_{\scriptstyle \pm 1.07}$ & 28.93$_{\scriptstyle \pm 0.43}$ & \cellcolor{gray!15}\textbf{20.10$_{\scriptstyle \pm 0.14}$} \\
 & \multirow{2}{*}{14B} & 0.5 & 11.43$_{\scriptstyle \pm 0.04}$ & 11.38$_{\scriptstyle \pm 0.04}$ & 10.79$_{\scriptstyle \pm 0.04}$ & \cellcolor{gray!15}\textbf{10.02$_{\scriptstyle \pm 0.03}$} & 11.51$_{\scriptstyle \pm 0.05}$ & 11.36$_{\scriptstyle \pm 0.02}$ & 10.92$_{\scriptstyle \pm 0.10}$ & \cellcolor{gray!15}\textbf{9.95$_{\scriptstyle \pm 0.13}$} \\
 &  & 0.7 & 37.04$_{\scriptstyle \pm 0.52}$ & 35.62$_{\scriptstyle \pm 0.27}$ & 34.29$_{\scriptstyle \pm 0.65}$ & \cellcolor{gray!15}\textbf{33.71$_{\scriptstyle \pm 0.45}$} & 20.82$_{\scriptstyle \pm 0.35}$ & 20.18$_{\scriptstyle \pm 0.16}$ & 20.37$_{\scriptstyle \pm 0.23}$ & \cellcolor{gray!15}\textbf{18.90$_{\scriptstyle \pm 0.28}$} \\
\bottomrule
\end{tabular}
\label{tab:ppl}
\vspace{-0.5cm}

\end{table*}

\begin{table*}[t]
\centering
\scriptsize
\setlength{\tabcolsep}{2.0pt}
\caption{\small Comparison of Zero-shot Accuracy  ($\uparrow$). We evaluate Base pruning methods (Wanda and SparseGPT) and their variants augmented with FARMS, ATP, and SWA at sparsity 0.5/0.7. Results are averaged over three seeds and reported as mean$_{\pm \text{std}}$. The best result in each setting is \colorbox{gray!15}{\textbf{bolded}}. SWA improves zero-shot accuracy and achieves the strongest performance over all the models and sparsity levels.}
\begin{tabular}{ccc|llll|llll}
\toprule
\multirow{2}{*}{\textbf{Model}} & \multirow{2}{*}{\textbf{Size}} & \multirow{2}{*}{\textbf{Sparsity}} & \multicolumn{4}{c|}{\textbf{Wanda}} & \multicolumn{4}{c}{\textbf{SparseGPT}} \\
 &  &  & Base & + FARMS & + ATP & \textbf{+ SWA} & Base & + FARMS & + ATP & \textbf{+ SWA} \\
\midrule
\multirow{4}{*}{\shortstack{Llama 3.2}} & \multirow{2}{*}{1B} & 0.5 & 34.91$_{\scriptstyle \pm 3.67}$ & 35.21$_{\scriptstyle \pm 2.28}$ & 37.16$_{\scriptstyle \pm 0.82}$ & \cellcolor{gray!15}\textbf{38.89$_{\scriptstyle \pm 1.03}$} & 39.57$_{\scriptstyle \pm 2.67}$ & 40.33$_{\scriptstyle \pm 2.41}$ & 41.26$_{\scriptstyle \pm 0.23}$ & \cellcolor{gray!15}\textbf{42.85$_{\scriptstyle \pm 1.75}$} \\
 &  & 0.7 & 25.86$_{\scriptstyle \pm 3.62}$ & 27.29$_{\scriptstyle \pm 4.75}$ & 29.86$_{\scriptstyle \pm 4.65}$ & \cellcolor{gray!15}\textbf{31.11$_{\scriptstyle \pm 4.81}$} & 30.85$_{\scriptstyle \pm 4.67}$ & 29.38$_{\scriptstyle \pm 1.33}$ & 35.29$_{\scriptstyle \pm 0.22}$ & \cellcolor{gray!15}\textbf{35.56$_{\scriptstyle \pm 0.34}$} \\
 & \multirow{2}{*}{3B} & 0.5 & 40.53$_{\scriptstyle \pm 1.41}$ & 42.71$_{\scriptstyle \pm 1.22}$ & 44.57$_{\scriptstyle \pm 1.39}$ & \cellcolor{gray!15}\textbf{44.72$_{\scriptstyle \pm 1.87}$} & 45.50$_{\scriptstyle \pm 0.36}$ & 46.67$_{\scriptstyle \pm 0.50}$ & 46.05$_{\scriptstyle \pm 1.41}$ & \cellcolor{gray!15}\textbf{48.49$_{\scriptstyle \pm 0.31}$} \\
 &  & 0.7 & 40.16$_{\scriptstyle \pm 2.41}$ & 41.50$_{\scriptstyle \pm 1.63}$ & 40.91$_{\scriptstyle \pm 4.74}$ & \cellcolor{gray!15}\textbf{43.04$_{\scriptstyle \pm 0.87}$} & 40.42$_{\scriptstyle \pm 1.16}$ & 42.25$_{\scriptstyle \pm 0.85}$ & 41.69$_{\scriptstyle \pm 4.46}$ & \cellcolor{gray!15}\textbf{43.57$_{\scriptstyle \pm 1.11}$} \\
\multirow{2}{*}{\shortstack{Llama 3}} & \multirow{2}{*}{8B} & 0.5 & 52.12$_{\scriptstyle \pm 1.25}$ & 53.49$_{\scriptstyle \pm 2.17}$ & 52.43$_{\scriptstyle \pm 1.32}$ & \cellcolor{gray!15}\textbf{55.22$_{\scriptstyle \pm 1.52}$} & 48.53$_{\scriptstyle \pm 2.24}$ & 50.28$_{\scriptstyle \pm 1.30}$ & 53.48$_{\scriptstyle \pm 0.61}$ & \cellcolor{gray!15}\textbf{54.56$_{\scriptstyle \pm 0.85}$} \\
 &  & 0.7 & 27.57$_{\scriptstyle \pm 0.07}$ & 27.14$_{\scriptstyle \pm 1.84}$ & 35.90$_{\scriptstyle \pm 0.74}$ & \cellcolor{gray!15}\textbf{36.75$_{\scriptstyle \pm 0.72}$} & 37.26$_{\scriptstyle \pm 5.39}$ & 38.38$_{\scriptstyle \pm 0.90}$ & 42.33$_{\scriptstyle \pm 2.86}$ & \cellcolor{gray!15}\textbf{43.32$_{\scriptstyle \pm 1.56}$} \\
\midrule
\multirow{6}{*}{\shortstack{Qwen 3}} & \multirow{2}{*}{4B} & 0.5 & 45.69$_{\scriptstyle \pm 3.35}$ & 46.33$_{\scriptstyle \pm 2.49}$ & 47.45$_{\scriptstyle \pm 0.39}$ & \cellcolor{gray!15}\textbf{47.85$_{\scriptstyle \pm 0.42}$} & 43.75$_{\scriptstyle \pm 1.25}$ & 47.28$_{\scriptstyle \pm 1.04}$ & 46.95$_{\scriptstyle \pm 4.51}$ & \cellcolor{gray!15}\textbf{48.00$_{\scriptstyle \pm 1.23}$} \\
 &  & 0.7 & 35.86$_{\scriptstyle \pm 0.06}$ & 36.60$_{\scriptstyle \pm 1.64}$ & 36.98$_{\scriptstyle \pm 0.48}$ & \cellcolor{gray!15}\textbf{38.00$_{\scriptstyle \pm 0.51}$} & 41.21$_{\scriptstyle \pm 4.51}$ & 42.50$_{\scriptstyle \pm 0.89}$ & 42.00$_{\scriptstyle \pm 4.02}$ & \cellcolor{gray!15}\textbf{43.02$_{\scriptstyle \pm 1.23}$} \\
 & \multirow{2}{*}{8B} & 0.5 & 52.50$_{\scriptstyle \pm 0.89}$ & 53.28$_{\scriptstyle \pm 0.77}$ & 55.38$_{\scriptstyle \pm 0.19}$ & \cellcolor{gray!15}\textbf{56.12$_{\scriptstyle \pm 0.24}$} & 53.04$_{\scriptstyle \pm 3.22}$ & 54.11$_{\scriptstyle \pm 2.70}$ & 55.64$_{\scriptstyle \pm 0.86}$ & \cellcolor{gray!15}\textbf{55.67$_{\scriptstyle \pm 0.66}$} \\
 &  & 0.7 & 38.95$_{\scriptstyle \pm 2.04}$ & 39.57$_{\scriptstyle \pm 0.53}$ & 39.64$_{\scriptstyle \pm 0.53}$ & \cellcolor{gray!15}\textbf{41.66$_{\scriptstyle \pm 0.88}$} & 42.24$_{\scriptstyle \pm 0.52}$ & 42.91$_{\scriptstyle \pm 1.72}$ & 43.48$_{\scriptstyle \pm 0.74}$ & \cellcolor{gray!15}\textbf{44.87$_{\scriptstyle \pm 0.67}$} \\
 & \multirow{2}{*}{14B} & 0.5 & 57.05$_{\scriptstyle \pm 0.44}$ & 57.31$_{\scriptstyle \pm 0.22}$ & 57.86$_{\scriptstyle \pm 1.15}$ & \cellcolor{gray!15}\textbf{58.02$_{\scriptstyle \pm 0.45}$} & 57.12$_{\scriptstyle \pm 0.39}$ & 57.62$_{\scriptstyle \pm 4.73}$ & 57.38$_{\scriptstyle \pm 0.15}$ & \cellcolor{gray!15}\textbf{58.10$_{\scriptstyle \pm 0.34}$} \\
 &  & 0.7 & 43.91$_{\scriptstyle \pm 3.11}$ & 44.69$_{\scriptstyle \pm 2.42}$ & 45.41$_{\scriptstyle \pm 0.93}$ & \cellcolor{gray!15}\textbf{46.44$_{\scriptstyle \pm 2.86}$} & 48.43$_{\scriptstyle \pm 2.22}$ & 49.11$_{\scriptstyle \pm 1.80}$ & 48.97$_{\scriptstyle \pm 2.01}$ & \cellcolor{gray!15}\textbf{49.93$_{\scriptstyle \pm 2.22}$} \\
\bottomrule
\end{tabular}
\label{tab:zero}
\end{table*}

\begin{table*}[t]
\centering
\scriptsize
\setlength{\tabcolsep}{6pt}

\caption{\small ARAOC task mismatch rate at sparsity 0.5. We report mismatch rate (Not M $\downarrow$) across six tasks. Results are averaged over five runs and reported as mean$_{\pm \text{std}}$. The best result in each setting is \colorbox{gray!15}{\textbf{bolded}}.}

\begin{tabular}{l|cccccc}
\toprule
\textbf{Method} 
& \textbf{Move} 
& \textbf{Change Color} 
& \textbf{Copy} 
& \textbf{Mirror} 
& \textbf{Fill Internal} 
& \textbf{Scale} \\
\midrule

Qwen3 8B (Non-Pruned) 
& 36.00$_{\scriptstyle \pm 4.11}$
& 23.14$_{\scriptstyle \pm 4.61}$
& 12.18$_{\scriptstyle \pm 4.12}$
& 84.66$_{\scriptstyle \pm 3.20}$
& 6.07$_{\scriptstyle \pm 1.99}$
& 83.33$_{\scriptstyle \pm 5.87}$ \\
\midrule

SparseGPT 
& 44.21$_{\scriptstyle \pm 9.20}$
& 48.87$_{\scriptstyle \pm 10.40}$
& 60.99$_{\scriptstyle \pm 10.67}$
& 94.76$_{\scriptstyle \pm 9.21}$
& 34.25$_{\scriptstyle \pm 4.67}$
& 88.91$_{\scriptstyle \pm 8.03}$ \\

SparseGPT + FARMS 
& 46.22$_{\scriptstyle \pm 8.21}$
& 50.77$_{\scriptstyle \pm 15.85}$
& 56.60$_{\scriptstyle \pm 12.80}$
& 96.66$_{\scriptstyle \pm 10.76}$
& 33.11$_{\scriptstyle \pm 5.11}$
& 88.06$_{\scriptstyle \pm 8.95}$ \\

SparseGPT + ATP 
& 43.08$_{\scriptstyle \pm 9.01}$
& 47.78$_{\scriptstyle \pm 8.66}$
& 56.77$_{\scriptstyle \pm 12.48}$
& 93.55$_{\scriptstyle \pm 9.63}$
& 32.65$_{\scriptstyle \pm 5.81}$
& 89.48$_{\scriptstyle \pm 7.47}$ \\

SparseGPT + \textbf{SWA} 
& \cellcolor{gray!15}\textbf{41.41$_{\scriptstyle \pm 9.39}$}
& \cellcolor{gray!15}\textbf{47.07$_{\scriptstyle \pm 9.68}$}
& \cellcolor{gray!15}\textbf{56.11$_{\scriptstyle \pm 11.62}$}
& \cellcolor{gray!15}\textbf{91.28$_{\scriptstyle \pm 8.99}$}
& \cellcolor{gray!15}\textbf{30.69$_{\scriptstyle \pm 4.17}$}
& \cellcolor{gray!15}\textbf{86.20$_{\scriptstyle \pm 8.17}$} \\

\bottomrule
\end{tabular}
\vspace{-0.5cm}

\label{tab:araoc_split}
\end{table*}
\begin{table}[t]
\centering
\small
\setlength{\tabcolsep}{6pt}
\caption{ \small Ablation study on hierarchical structural deletion. We report the absolute WikiText perplexity (PPL ↓) of SWA and the PPL increase of each ablated variant relative to SWA under sparsity 0.5 and 0.7. SWA(Layer) uses only layer-level scores, and SWA(Head) uses only head-level scores.}
\label{tab:ablation_hier}
\begin{tabular}{lcccc}
\toprule
\multirow{2}{*}{\textbf{Variant}} 
& \multicolumn{2}{c}{\textbf{SparseGPT}} 
& \multicolumn{2}{c}{\textbf{Wanda}} \\
\cmidrule(lr){2-3} \cmidrule(lr){4-5}
& \textbf{0.5} & \textbf{0.7} 
& \textbf{0.5} & \textbf{0.7} \\
\midrule
\textbf{SWA} 
& \textbf{11.27} & \textbf{49.56} 
& \textbf{13.32} & \textbf{120.65} \\
SWA(Layer) 
& $+3.31$ & $+11.37$ 
& $+0.05$ & $+6.96$ \\
SWA(Head) 
& $+3.63$ & $+19.56$ 
& $+0.24$ & $+46.51$ \\
\bottomrule
\end{tabular}
\vspace{-0.5cm}
\end{table}



We conduct experiments to evaluate whether small-world functional organization is associated with LLM reasoning ability, and whether scores derived from this organization can preserve small-world structure during pruning. Our evaluation is guided by three research questions: \textbf{RQ1.} Does SWA preserve small-world functional organization when used to guide pruning? \textbf{RQ2.} Does pruning with SWA better maintain model performance under sparsity than existing pruning baselines? \textbf{RQ3.} Does pruning with SWA preserve intelligence-related reasoning ability on held-out reasoning benchmarks?

\vspace{-0.2cm}
\subsection{Experimental Setup}
\label{sec:swa_exp_setup}

First, we describe the setup used to collect activations, construct attention-head functional graphs, and compute SWI. Second, we describe the pruning setup used to evaluate whether SWA preserves performance and reasoning ability under sparsity. Details are provided in Section~\ref{app:swa_expdetails}.


\paragraph{Setup for SWI analysis.}
We first evaluate the relationship between SWI and reasoning performance using six open-source LLMs: Qwen3-4B, Qwen3-8B, Qwen3-14B, Llama-3.2-1B, Llama-3.2-3B, and Meta-Llama-3-8B \cite{yang2025qwen3, grattafiori2024llama}. We collect attention-head activations on DRE-Bench \cite{yang2025truly}, using Level 1 and Level 2 questions as calibration inputs for constructing attention-head functional graphs. We exclude Level 3 and Level 4 because current open-source models often fail on these harder tasks, making their generated answers unreliable for activation-based analysis. We evaluate DRE-Bench performance by computing the mean per-token negative log-likelihood over the answer tokens (mean answer NLL/token; lower is better) and compare it with each model's SWI.

\paragraph{Setup for pruning evaluation.}
For pruning experiments, we use the same six LLMs and apply SWA to SparseGPT \cite{frantar2023sparsegpt} and Wanda \cite{sun2023simple}. We compare against uniform sparsity and layer-wise allocation baselines, including FARMS \cite{hu2025eigenspectrum} and ATP \cite{huang2025determining}. We evaluate pruned models using WikiText perplexity \cite{merity2016pointer} and average zero-shot accuracy over BoolQ \cite{clark2019boolq}, RTE \cite{wang2018glue}, HellaSwag \cite{zellers2019hellaswag}, WinoGrande \cite{sakaguchi2021winogrande}, ARC-Easy, ARC-Challenge \cite{clark2018think}, and OpenBookQA \cite{mihaylov2018can} datasets.


\subsection{Effectiveness of Preserving SWI}

To answer \textbf{RQ1}, we evaluate whether SWA preserves small-world functional organization during pruning. For each of the six models, we apply SparseGPT at a sparsity level of $0.5$ and compare it with variants using our score-guided allocation strategies. We then recompute the SWI of the resulting head--head functional graphs. \textbf{SWA (Head)} and \textbf{SWA (Layer)} denote head-only and layer-only allocation, respectively. As shown in Table~\ref{tab:swi_pruning_preservation}, SWA consistently achieves higher SWI values than the baseline pruning methods. Although the absolute SWI decreases substantially after pruning, SWA retains higher relative SWI than competing allocation strategies under the same sparsity level.





\subsection{Effectiveness of SWA}
To answer \textbf{RQ2}, we evaluate the performance of SWA on the WikiText validation set and report the perplexity (PPL) in Table \ref{tab:ppl}. We find that our method consistently outperforms the baselines across all models and sparsity levels, demonstrating that preserving small-world functional organization during pruning leads to better model performance. Notably, at higher sparsity levels, the performance gap between our method and the baselines becomes more pronounced, indicating that SWI preserved pruning is particularly beneficial when aggressive pruning is applied. Additionally, we report the zero-shot performance on the downstream tasks in Table \ref{tab:zero}. Our method also achieves superior performance compared to the baselines, further confirming that preserving small-world functional organization during pruning helps maintain the model's reasoning abilities on intelligence benchmarks.

\subsection{Effectiveness of SWA in Preserving Intelligence-Related Reasoning Performance}

To further evaluate the effectiveness of SWA in preserving intelligence-related reasoning ability (\textbf{RQ3}), we conduct experiments on the ARAOC benchmark \cite{wu2025understanding}, which consists of a diverse set of reasoning tasks designed to assess various aspects of intelligence. We report the performance of our method and the baselines on the ARAOC benchmark in Table \ref{tab:araoc_split}. Our method consistently outperforms the baselines across all tasks, demonstrating that structure-aware pruning not only preserves small-world functional organization but also helps maintain the model's reasoning abilities on a wide range of intelligence-related tasks. This further validates the importance of considering structural properties when determining sparsity rates for pruning LLMs.

\subsection{Ablation Study}

We conduct an ablation study to examine the contribution of each component of SWA. To isolate the effect of each component, we compare it with two ablated variants: \textbf{SWA(Layer)}, which uses the layer-only scores, and \textbf{SWA(Head)}, which uses the head-only scores. We evaluate these variants with both SparseGPT and Wanda as the base pruning methods, using WikiText perplexity under sparsity levels 0.5 and 0.7. As shown in Table~\ref{tab:ablation_hier}, removing either component leads to higher perplexity, indicating that both head-level and layer-level structural signals contribute to preserving model quality after pruning. The degradation is especially pronounced at higher sparsity. For SparseGPT, removing the layer-level component increases perplexity by 19.56 at sparsity 0.7, while removing the head-level component increases perplexity by 11.37. A similar pattern holds for Wanda: at sparsity 0.7, removing the layer-level component causes a larger drop than removing the head-level component. These results suggest that layer-level scores provide the main allocation signal, while head-level scores refine pruning within layers. Overall, the ablation supports the hierarchical design, especially under high sparsity.
